%% kapa: exam=1 round=25 year=2014 period=2 form=A subject=accounting subject_ko=회계학 seq=6 number=46
%% source: Q-Net 감정평가사 기출문제 (cst00309) artlSeq=5101207 file=extracted/1008/5101207/2014년 제25회 감정평가사 1차시험 기출문제/2교시(A형).hwp
%% conversion: kordoc 4.24.51 -> tools/gen1.py (auto); review=proofread-sonnet; flags=none
\begin{question}{46}
재무제표 표시에 관한 설명으로 옳은 것은?
\begin{choices}
  \item 부적절한 회계정책은 이에 대하여 공시나 주석 또는 보충 자료를 통해 설명함으로써 정당화될 수 있다.
  \item 비유동자산의 처분손익을 처분대금에서 그 자산의 장부금액과 관련처분비용을 차감하여 표시하는 것은 총액주의에 위배되므로 허용되지 아니한다.
  \item 재무제표 항목의 표시와 분류는 한국채택국제회계기준에서 표시방법의 변경을 요구하는 경우 이외에는 매기 동일하여야 한다.
  \item 기업이 기존의 대출계약조건에 따라 보고기간 후 적어도 12개월 이상 부채를 차환하거나 연장할 것으로 기대하고 있고, 그런 재량권이 있다하더라도, 보고기간 후 12개월 이내에 만기가 도래한다면 유동부채로 분류한다.
  \item 단기매매 금융상품에서 발생하는 손익과 같이 유사한 거래의 집합에서 발생하는 차익과 차손은 순액으로 표시한다. 그러나 그러한 차익과 차손이 중요한 경우에는 구분하여 표시한다.
\end{choices}
\end{question}
\answer{5}
